\chapter{Introduction and Thesis Claims}
\label{ch:introduction-thesis-claims}

A battery controller can be shown stale feeder telemetry and continue dispatching as if the
network were quiet. A vehicle can track a nominal lane geometry after the observation stream
has already fallen behind the scene. A bedside monitoring system can smooth away a sudden
physiological transition while its downstream policy still appears internally consistent. In
each case the failure is structurally the same: the controller acts on an observation
surface that is coherent enough for computation but no longer trustworthy enough for safe
release. This dissertation begins from that release-boundary hazard.

\section{From degraded observation to unsafe release}
Physical AI systems always act through an observation layer. Telemetry arrives late, state
estimators compress uncertainty, learned models summarize partially observed trajectories,
and communication channels drop, reorder, or distort critical signals. Most stacks treat
those effects as nuisances to be absorbed by better prediction, better perception, or better
control tuning. ORIUS starts from the opposite premise: once degraded observation changes
what can be known about the physical world, it changes what actions are safe to release.

The central danger is therefore not forecast error alone and not perception error alone. It
is the possibility that action legality on the observed state diverges from action legality
on the true state. ORIUS names that discrepancy the
\emph{observation--action safety gap} (OASG). The OASG is the dissertation's scientific
object because it identifies the exact place where nominal intelligence can remain coherent
while physical safety has already been lost.

\section{Thesis claim}
This dissertation advances one central claim: \emph{physical AI under degraded observation
requires a runtime safety layer that sits between nominal intelligence and physical release,
and that layer can be expressed through one reusable grammar of detection, uncertainty
inflation, safe-set tightening, action repair or fallback, and certificate-aware release.}
ORIUS is proposed as that layer.

The claim is strong but intentionally bounded. ORIUS is not a universal optimal controller,
not a replacement for domain-specific control design, and not yet a proof of equal-domain
closure across every supported application. It is a runtime discipline for deciding what
action may still be released once observation quality has degraded.

\section{ORIUS as the dissertation object}
The reusable object in this dissertation is a typed
Detect--Calibrate--Constrain--Shield--Certify kernel.
\begin{itemize}
  \item Detect scores reliability loss in the observation channel.
  \item Calibrate widens the observation-consistent state set.
  \item Constrain recomputes admissible action under that widened set.
  \item Shield repairs or replaces unsafe candidate actions.
  \item Certify records why the release remains admissible and how long that claim should
  remain valid.
\end{itemize}
The important point is that ORIUS does not compete with every controller on the same terms.
It addresses the release-boundary decision that appears once a controller admits the
observation channel may already be wrong.

\begin{figure}[htbp]
\centering
\includegraphics[width=0.9\textwidth]{fig_universal_framework}
\caption{ORIUS as a universal runtime safety layer between nominal intelligence and
physical release. The same five-stage kernel recurs across domains even though the domains
do not share a plant model, controller family, or optimization objective.}
\label{fig:intro-universal-framework}
\end{figure}

Three implementation-grounded surfaces make this claim more than conceptual language.
\texttt{DomainAdapter} exposes the domain safety predicate, true-state margin, and fallback
semantics that the runtime must protect. The universal benchmark schema exposes one replay
contract for true-state violation, intervention, fallback, certificate validity, and
latency. CertOS supplies the lifecycle semantics that determine when repaired action can be
issued, renewed, expired, revoked, or replaced by fallback.

Those three surfaces are not supporting details around a more abstract contribution. They
are the mechanism that makes ORIUS scientifically concrete. Without \texttt{DomainAdapter},
the safety object stays implicit and domain-local. Without the universal benchmark schema,
cross-domain degraded-observation claims remain incomparable. Without CertOS, repaired
action lacks a governed release semantics. Taken together, they are what turns ORIUS from a
collection of useful ideas into a reusable runtime layer.

\section{Main contributions}
The dissertation makes exactly four contributions.
\begin{enumerate}
  \item \textbf{A runtime object.} It defines OASG-centered safety under degraded
  observation as a first-class runtime problem for physical AI.
  \item \textbf{A systems object.} It presents the ORIUS kernel as an adapterized runtime
  safety layer built around \texttt{DomainAdapter}, the universal benchmark schema, and
  CertOS release semantics.
  \item \textbf{A formal object.} It develops a theorem bridge from reliability-conditioned
  uncertainty and repaired action to bounded violation reduction under explicit assumptions.
  \item \textbf{An empirical object.} It assembles bounded multi-domain evidence with one
  witness row and explicit defended and gated rows rather than flattening all domains into a
  false claim of equal maturity.
\end{enumerate}

\section{Claim discipline}
The manuscript uses three claim classes and keeps them distinct.
\begin{description}
  \item[Proved claims] theorem-level results under stated assumptions for OASG, repair, and
  bounded safety improvement.
  \item[Validated claims] replay-backed and runtime-governed empirical claims tied to
  tracked artifacts.
  \item[Roadmap claims] extensions that are structurally motivated but not promoted as
  defended evidence.
\end{description}
That distinction keeps the dissertation from confusing architectural universality with
empirical parity. The exact row-by-row maturity boundary is stated in the cross-domain
chapter and then restated once in the conclusion, rather than repeated throughout the
manuscript.

\section{Why ORIUS matters beyond one domain}
The importance of ORIUS is not that batteries, vehicles, industrial plants, healthcare
systems, navigation stacks, and aerospace controllers secretly reduce to the same plant. The
importance is that they all face the same release-boundary question once observation trust
degrades: what action can still be justified, with what margin, and under what fallback and
certificate? ORIUS freezes the semantics of that question so that cross-domain comparison
happens at the level of runtime safety rather than only at the level of application area.

\section{Dissertation organization}
The rest of the dissertation follows that argument directly. Chapters~\ref{ch:related-work-novelty-gap}
and~\ref{ch:formal-problem-formulation} define the novelty gap and formal problem object.
Chapters~\ref{ch:orius-architecture-dc3s} through~\ref{ch:benchmark-reproducibility-protocol}
present the ORIUS kernel, theorem ladder, and replay protocol. Chapters~\ref{ch:witness-domain-implementation}
and~\ref{ch:witness-results-failure-analysis} defend the battery witness row. Chapters~\ref{ch:universal-orius-across-domains}
through~\ref{ch:certos-runtime-assurance} test how far the same runtime grammar carries
across domains and governance surfaces. Chapter~\ref{ch:governance-reproducibility-limitations-conclusion}
then states the precise evidence boundary and the final thesis claim.
